\section{Distribution consequences}
\label{sec:consequences}

\subsection{Local counts and multiplicities}

All zeros in the slab $0<\Re s<1$, $T\le\Im s\le T+1$ lie in the disk
of radius $5/2$ around $2+i(T+1/2)$.  Apply (J) with outer radius four and
inner radius $5/2$.  Approximating $\zeta$ by a rational polynomial closely
enough to preserve any supplied finite family of zero-cluster certificates
gives
\[
 \sum_j \operatorname{count}(\mathcal B_j)
 \le\frac{\log(16M)}{\log(8/5)}
 \le12\log(T+11)+12,\qquad T\ge8.                       \tag{LC}
\]
Here the $\mathcal B_j$ are disjoint certified clusters in the slab and
$M=4(\lceil T\rceil+10)^4$.
The inequalities $\log(8/5)\ge3/8$ and $\log64<9/2$ suffice for the last
estimate.  This proves the usual $O(\log T)$ local zero-count bound and the
same bound on the multiplicity of an individual zero.  The certificate
form does not require a decision procedure for exact equality of ordinates.

\subsection{Infinitely many zeros, with a bounded witness}

Take an integer $T\ge2^{16}$ and choose regular heights
$U\in(T,T+1)$ and $V\in(2T,2T+1)$.  For $x\in[T+1,2T]$,
\[
 \mathcal M'(x)=\frac1{2\pi}\log\frac{x}{2\pi}
       \ge\frac1{16}\log T,
\]
since $2\pi<8$, $\log(2\pi)<3$, and $\log T\ge8$.
Thus $\mathcal M(V)-\mathcal M(U)\ge T\log T/32$.
The sum of the two errors in (RVM) is at most $30\log T$.
As $T\ge2^{16}>1920$,
\[
 N(V)-N(U)\ge\frac{T\log T}{64}\ge\frac T8.            \tag{AB}
\]
The two contour certificates therefore give a positive-count rectangle
strictly above height $T$, and indeed a count at least $T/8$.
Using $T=4^{m+8}$ gives disjoint such height ranges for successive $m$.
This is a bounded constructive witness for unboundedly many nontrivial
zero clusters in disjoint regions. The following construction extracts
a primitive recursive point name from each such cluster, yielding
a uniform sequence of pairwise apart nontrivial zeros. It does not
order all exact zero ordinates.

\subsection{Extracting an actual zero name from a positive count}
\label{sec:zero-extraction}

A positive contour count contains more constructive information than an
unrelated sequence of small residuals. Here is a coherent root extraction,
using only the quantitative minimum-modulus construction and the finite
winding rules already proved.

Fix one rational rectangle $Q_0$ with a certified zero-free boundary and
positive xi winding. Choose an integer $A\ge1$ so that
$|s-2|\le A$ throughout $Q_0$, and use $f(z)=\xi(2+Az)$. The entire
theta-integral construction gives, with $R=4A+2$,
\[
        |f(z)|\le M:=1+2R(R+1)R!\quad(|z|\le4),
             \qquad f(0)=\xi(2)=\zeta(2)/\pi\ge1/4.
\]
The same expression with $R=5A+2$ bounds the larger disk used for Taylor
approximation. The minimum-modulus lemma therefore applies with the fixed
integer $K=4\lceil\log_2(16M)\rceil$. In its notation one may take
\[
       L_\delta=2^{-83}M^{-20}(\delta/6)^K,
\]
so that outside at most $K$ disks of radius $\delta$ the normalized
function has modulus at least $3L_\delta/4$. The global derivative on
$Q_0$ is bounded by a fixed primitive recursive integer $D\ge1$; for
example the Cauchy bound from the displayed disk estimate permits $D=M$.
These constants may be very large, but none is obtained by an unbounded
search.

For a current rectangle $[a,a+w]\times[b,b+h]$, use the $K+1$ candidates
\[
 x_j=a+w/3+\frac{(j+1)w}{3(K+2)},\qquad
 y_j=b+h/3+\frac{(j+1)h}{3(K+2)}\quad(0\le j\le K),
\]
and put $\delta=\min(w,h)/(9A(K+2))$. Twice this radius is smaller
than either candidate spacing in normalized coordinates, so each exceptional
disk excludes at most one vertical candidate and at most one horizontal
candidate. There is therefore a safe full segment in each direction.
On every candidate segment evaluate $\xi$ to error $L_\delta/16$ on a
mesh of physical width at most $L_\delta/(8D)$. Choose the first segment
whose rational samples have modulus at least $L_\delta/2$, tested by
squaring. A safe candidate passes this test, and every passing candidate
has true modulus at least $5L_\delta/16$ on the entire segment. Record
the smaller margin $L_\delta/4$ for each of the two chosen cuts.

If the parent's recorded margin is $m$, give every child the common margin
\[
                         m'=\tfrac12\min(m,L_\delta/4).
\]
Every child edge then has true
modulus at least $2m'$. Restrict the inherited boundary certificates
to the child edges, and refine both inherited edges and the two cuts to
a common mesh of width at most $m'/(8D)$, with sample error at most $m'/8$.
The resulting Lipschitz certificates verify the claimed margin $m'$.
Reuse
identical samples with opposite orientations on shared child edges. This
constructs full closed-boundary certificates for all four children, not
just certificates for their new cuts. Each side length lies between
$1/3$ and $2/3$ of the corresponding parent side length.

Each child winding is nonnegative. To see this without already assuming
analytic root existence, approximate $\xi$ on that child's boundary by a
rational complex polynomial to less than half the certified boundary
margin. The finite zero-free homotopy identifies the two windings. The
polynomial engine of Section~\ref{sec:polynomial-algebra} factors that
polynomial, so its winding is the sum of the nonnegative indices of its
linear factors. Cancellation of the new internal edges makes the four
child windings sum to the parent's winding. At least one is consequently
positive; choose the first such child in a fixed ordering.

Iterating this procedure gives nested rational rectangles $Q_j$ with
positive winding and diameter at most $D_0(2/3)^j$, for a supplied
$D_0\ge1$. All searches are bounded, and the bounds are recomputed by
primitive recursion from the current rational rectangle and its inherited
boundary margin. In particular this is not a search that waits until a
possibly vanishing sample becomes visibly nonzero.

Let $c_j$ be the centre of $Q_j$. Then
\[
                        |\xi(c_j)|\le D\,\operatorname{diam}(Q_j).
\]
Indeed, a strict violation would put the entire boundary image in a disk
centred at $\xi(c_j)$ which misses zero; its winding would be zero. The
finite-precision version proves the displayed weak inequality by ruling
out every rationally witnessed violation, using the same derivative bound.

Choose an integer $L$ with $2^L\ge4DD_0$, and set
\[
                         \alpha(k)=c_{\,2(k+L)}.
\]
Since $(2/3)^2<1/2$, nestedness gives, for $\ell\ge k$,
\[
       |\alpha(\ell)-\alpha(k)|\le2^{-k-2},\qquad
       |\xi(\alpha(k))|\le2^{-k-2}.
\]
The centres therefore form one primitive recursive fast Cauchy name for an
actual zero. The initial positive boundary margin and derivative bound
also keep this zero positively separated from $\partial Q_0$, so it lies
inside the original rectangle. Every recurrence invariant used to obtain these inequalities
concerns finite rational rectangles, contour certificates, signed integer
counts and fixed-precision approximation errors. No countable choice or
decision of exact zero ordinates is involved.

For rectangles above the real axis, the explicitly nonvanishing factors
relating $\xi$ and $\zeta$ identify the extracted point as a nontrivial
zeta zero. Applying the construction to disjoint certified height bands
gives pairwise apart zero names. This proves constructive infinitude when
the counting estimate makes those band counts positive; it does not supply
a procedure for deciding equal ordinates or sorting all exact zeros.


\subsection{Average density in growing windows}

From (RVM), for $T\le U<T+1$ and
$T+H\le V<T+H+1$, where $1\le H\le T$, one obtains
\[
 N(V)-N(U)=\frac{H}{2\pi}\log\frac T{2\pi}
          +O\left(\frac{H^2}{T}+\log(T+H+4)\right).
\]
The constants here can be read directly from (RVM) and the derivative of
$\mathcal M$; alternatively one retains the exact difference $\mathcal M(V)-\mathcal M(U)$ and the
explicit error sum.  Hence when $H\to\infty$ and $H=o(T)$, the mean density
in such a window is asymptotic to $\log T/(2\pi)$.
For a unit window the error is the same order as the main term; the correct
consequence of this proof is only $N(T+1)-N(T)=O(\log T)$, not an
asymptotic equivalent.  Similarly, the usual inverse relation
$\gamma_n\sim2\pi n/\log n$ is a classical consequence of the count
asymptotic.  Its constructive reading is via inverse count brackets and
does not silently posit a uniformly decidable ordering of equal zero
ordinates.

\subsection{An explicit growing-window estimate}
The asymptotic just stated can be given a convenient numerical form.
Let $T\ge8$, $1\le H\le T$, and choose the certified heights
$U\in(T,T+1)$ and $V\in(T+H,T+H+1)$. By
\eqref{eq:main-rvm},
\begin{equation}
 \left|N(V)-N(U)-\frac{H}{2\pi}\log\frac{T}{2\pi}\right|
 \le\frac{H^2}{12T}+18\log(T+H+5)+80.
 \label{eq:window-explicit}
\end{equation}
Indeed the change from $U,V$ to $T,T+H$ in the main-term difference
is at most $2\log(T+H+5)$. The change of the logarithm inside
$\int_T^{T+H}\mathcal M'(x)dx$ contributes at most
$H^2/(4\pi T)\le H^2/(12T)$. Finally the two counting errors are
at most $16\log(T+H+5)+80$.

For a particularly simple arithmetic specialization take $T=n^2$,
$H=n$, with integer $n\ge3$. The right side becomes
$1/12+18\log(n^2+n+5)+80$, while the leading term is
$n\log(n^2/(2\pi))/(2\pi)$. Dividing proves the claimed mean density
with a primitive recursive convergence modulus. One need not introduce
an irrational square-root input into the checker.

\subsection{Effective tails of reciprocal-zero sums}
\label{sec:zero-sums}

For every certified height $u\ge8$, the coarse consequence
\begin{equation}
                    N(u)\le64u\log(u+2)
 \label{eq:coarse-count}
\end{equation}
follows immediately from~\eqref{eq:main-rvm}. For example, the main
term is at most $u\log(u+2)/6$, and its remaining error is at most
$36\log(u+2)\le(9u/2)\log(u+2)$. The factor $64$ leaves ample
slack for rational endpoint adjustments.

\begin{theorem}[Uniform reciprocal-square tail]
For every integer $T\ge8$, every finite collection of zeros with
$|\Im\rho|>T$, counted with multiplicity, satisfies
\begin{equation}
       \sum |\rho|^{-2}
       \le\frac{2048(\log(T+2)+1)}{T}.
 \label{eq:reciprocal-tail}
\end{equation}
The finite constructive formulation uses disjoint certified clusters
and the lower bound $|\rho|\ge|\Im\rho|$ in each height band.
\end{theorem}

\begin{proof}
Partition positive ordinates into the dyadic bands
$(2^jT,2^{j+1}T]$. To include a possible zero at the upper endpoint,
choose a certified height $u_j\in(2^{j+1}T,2^{j+1}T+1)$.
By~\eqref{eq:coarse-count}, the count of positive ordinates in that
band is at most
\[
 144\,2^jT\bigl(\log(T+2)+(j+1)\log2\bigr).
\]
Here $2^{j+1}T+1\le(9/8)2^{j+1}T$ and
$2^{j+1}T+3\le2^{j+1}(T+2)$ suffice. Conjugation doubles this
bound for both signs. Division by $(2^jT)^2$ and summation give
\[
 \frac{288}{T}\sum_{j\ge0}2^{-j}
       \bigl(\log(T+2)+(j+1)\log2\bigr)
 =\frac{576\log(T+2)+1152\log2}{T},
\]
which is smaller than~\eqref{eq:reciprocal-tail}.
The actual arithmetic argument uses a finite geometric sum and its
explicit remainder. It proves the same bound for every finite number
of bands, so no induction over a quantified convergence assertion is
being used.
\end{proof}

It follows that $\sum_\rho\rho^{-m}$ is absolutely convergent for
each integer $m\ge2$. For $m=1$, pair each positive-height zero
$\rho$ with the negative-height zero $1-\rho$. The identity
\begin{equation}
 \frac1\rho+\frac1{1-\rho}
       =\frac1{\rho(1-\rho)}
 \label{eq:paired-reciprocals}
\end{equation}
gives modulus at most $1/(\Im\rho)^2$. Hence this paired reciprocal
series is absolutely convergent as a series of pairs, without a
Riemann-hypothesis assumption. This is not a claim of absolute
convergence of the unpaired series $\sum_\rho1/\rho$.

An explicit uniform tail choice for either assertion is
\begin{equation}
                         T_p=2^{2p+32}.
 \label{eq:reciprocal-modulus}
\end{equation}
Since $\log(T_p+2)+1\le2p+34$ and
$p+17\le2^{p+5}$, the right side of
\eqref{eq:reciprocal-tail} is less than $2^{-p}$.

The finite partial sums also have a constructive representation which
does not list exact zeros. Put $\eta=1/4096$. The derivative bound for
$\xi$ with $R=3$ is $D_3=372$, so its positive real-axis bound gives
$|\xi(x+\ii\eta)|>1/4$ for $-1\le x\le2$.
On the certified rectangle $[-1,2]+\ii[\eta,U]$, integrate
\[
 \frac{1}{2\pi\ii}\frac{\xi'(s)}{\xi(s)}w(s),
 \quad w(s)=s^{-m}\quad\hbox{or}\quad w(s)=\frac1{s(1-s)}.
\]
The weights have no pole in this rectangle. Division has a positive
boundary clearance; derivatives and quadrature have the explicit bounds
already established. Polynomial factorization and the elementary
one-factor residue calculation identify this integral with the
corresponding finite zero sum, counting multiplicity. For $m\ge2$
include the reflected lower rectangle as well; for the paired sum the
upper rectangle alone suffices. Approximate these contour integrals at
precision $p+2$ and use~\eqref{eq:reciprocal-modulus} for the tail.
This supplies primitive recursive Cauchy names for the stated sums.
